\subsection{The entropy comparison in the macroscopic range}
The comparison in this subsection reveals a placement in uniform order and bounds, by Jensen's
inequality, the logarithm of the available mass at each step. This is the scheme of the entropy
method of Radhakrishnan~\cite{radhakrishnan} and of Linial and Luria~\cite{linialluria}, which
Luria~\cite{luria} applied to the upper bounds for the queens problem.

\begin{lema}[Factorization at a stationary point]\label{coste-vacancias}
Let $p$ be an interior stationary point of the Bethe entropy $H_B$, with row and column masses
equal to one and $h_f>0$. There are positive numbers $a_i,b_j$ such that, on every available square,
\[
 v_e:=p_e(1-p_e)^3=a_i b_j h_{D(e)}h_{A(e)} .
\]
Set $\mathcal K=\prod_i a_i\prod_j b_j\prod_f h_f$ and $W=\sum_e v_e$. Then
$H_B(p)=-\log\mathcal K+3\sum_e\log(1-p_e)$. If also $p_{\max}\le C/k$, then for
sufficiently large $k$ we have $W=k+O_C(1)$ and $H_B(p)=-\log\mathcal K-3k+O_C(1)$.
\end{lema}
\begin{proof}
The derivative of $H_B$ with respect to $p_e$ is
$-\log p_e-3\log(1-p_e)+\log h_{D(e)}+\log h_{A(e)}-2$.
At the interior stationary point, the gradient has zero inner product with every variation
preserving the row and column sums. By orthogonality to the kernel of the map giving these sums,
the derivative is the sum of a row multiplier and a column multiplier. Exponentiating and
absorbing the constant gives the factorization. Taking logarithms, multiplying by $p_e$, and
summing over the squares, we obtain $\sum_i\log a_i+\sum_j\log b_j$ from the two sets of
multipliers and $\sum_f\ell_f\log h_f$ from the diagonals.
The identity for $H_B$ follows from $\ell_f+h_f=1$. Moreover, $0\le p_e-v_e\le3p_e^2$ and
$|\log(1-p_e)+p_e|\le p_e^2/(1-p_{\max})$. Both bounds can be summed because $\sum_e p_e^2\le C$.
\end{proof}

\begin{lema}[Comparison with augmented placements]\label{coste-comparacion}
Let $B$ have $Q(B)>0$ and an interior stationary Bethe point $p$ with unit row and column
masses and positive diagonal vacancies. Use the real-piece weights $v_e=p_e(1-p_e)^3$ and the
product $\mathcal K$ from the factorization in Lemma~\ref{coste-vacancias}.
With $n$ the order of the initial board, put $N=2n-1$ and $H=N-k$.
There is a comparison law, uniform on ordered augmented placements with $k$ real pieces and $H$
virtual pieces, for which
\[
 \log Q(B)\le
 \E\sum_{s=1}^N\log W_{T_s}(s)
       -\log\mathcal K-\log k!-4\log H! .
\]
Here $T_s$ is the type of the piece in position $s$, and $W_{T_s}(s)$ is the available mass
of that type just before the piece is revealed.
\end{lema}
\begin{proof}
We adapt the count of augmented placements in \cite[Sections~3.1--3.4 and
Theorem~3.6]{ultimas} to the support of $B$.
Take $N$ vertices in each of the four families. In the first two families, retain the $k$ physical
rows and $k$ physical columns of $B$ and add $H$ dummy vertices to each family. In each diagonal
family, use the $N$ labels of the initial board. A residual diagonal retains its physical line.
In place of a diagonal already occupied, put a dummy vertex with the same label, incident to no
residual square. A line with no available squares also has real degree zero. In both cases the
residual mass is zero and we extend the vacancy to this label by setting $h_f=1$. This extension enters
only the auxiliary count. It preserves the stationarity identity on every available square and
the value of $\mathcal K$.

The real pieces are the available squares, with weights $v_e$. For the virtual pieces we assign
vertex weight zero to physical rows and columns, one to dummy rows and columns, and $h_f$ to
diagonal labels. A virtual piece consists of one vertex from each family and has weight equal to
the product of the four vertex weights. We use only pieces with positive weight. In a diagonal
family each square contributes its mass exactly once. Hence
\[
 \sum_{f\text{ in one diagonal family}}h_f=N-\sum_e p_e=H.
\]
The vertex weights sum to $H$ in each of the other two families as well. Thus the initial virtual
mass is $H^4$, even if some diagonals were occupied before reaching $B$.

A complete real placement leaves $H$ vertices uncovered in each family. To complete it, use the
$H$ dummy rows as indices and independently choose three bijections from them to the uncovered
vertices of the other three families. This gives exactly $(H!)^3$ virtual completions, all of
positive weight since every $h_f>0$. Conversely, in an augmented placement covering every vertex,
real pieces must cover the physical rows and virtual pieces must cover the dummy rows. The real
pieces thus form one of the $Q(B)$ completions of $B$. Each augmented placement admits $N!$
orderings. In every ordering the product of the weights of all the pieces is $\mathcal K$.
Indeed, the real pieces contribute each $a_i$ and $b_j$ once, and the real and virtual pieces
together use the factor $h_f$ of each diagonal label once.

As in~\cite[Section~3, Definitions~3.1 and~3.2]{ultimas}, let $\mathbb P$ be the comparison law,
uniform on these successful
histories, and let $\mathbb Q$ be the following reference process. First choose a uniformly random
color word with $k$ real and $H$ virtual colors. At each position choose an available piece of the
prescribed type with probability proportional to its weight. A piece is available if all its
vertices are still uncovered. If the mass of that type is zero, the process terminates in the
failure state, which retains the corresponding probability. For a successful history $\omega$,
\[
 \mathbb P(\omega)=\frac{1}{Q(B)(H!)^3N!},\qquad
 \mathbb Q(\omega)=\frac{k!H!}{N!}\,
             \frac{\mathcal K}{\prod_{s=1}^N W_{T_s}(s)} .
\]
In particular, $\mathbb Q(\omega)>0$. Jensen's inequality for the logarithm gives
\[
 \E_{\mathbb P}\log(\mathbb Q/\mathbb P)\le\log\sum_{\omega:\mathbb P(\omega)>0}\mathbb Q(\omega)\le0 .
\]
Substituting the two probabilities gives the asserted inequality. The expectation is under
$\mathbb P$ on the single board $B$ whose completions we are counting.
\end{proof}
